\documentclass[reqno]{amsart} 
\usepackage{amssymb,latexsym,amsmath,amscd,graphicx,setspace,amsthm,verbatim,enumitem}
\usepackage[margin = 3 cm]{geometry}

\def\upint{\mathchoice%
    {\mkern13mu\overline{\vphantom{\intop}\mkern7mu}\mkern-20mu}%
    {\mkern7mu\overline{\vphantom{\intop}\mkern7mu}\mkern-14mu}%
    {\mkern7mu\overline{\vphantom{\intop}\mkern7mu}\mkern-14mu}%
    {\mkern7mu\overline{\vphantom{\intop}\mkern7mu}\mkern-14mu}%
  \int}
\def\lowint{\mkern3mu\underline{\vphantom{\intop}\mkern7mu}\mkern-10mu\int}


\theoremstyle{plain}
\newtheorem{theorem}{Theorem}[section]
\newtheorem{proposition}{Proposition}
\newtheorem{corollary}{Corollary}
\newtheorem{lemma}{Lemma}
\newtheorem{conjecture}{Conjecture}
\newtheorem{question}{Question}
\newtheorem{problem}{Problem}
      
\theoremstyle{definition}
\newtheorem{definition}{Definition}

\newenvironment{solution1}{\paragraph{\emph{Solution $1$}.}}{\hfill$\square$}
\newenvironment{solution2}{\paragraph{\emph{Solution $2$}.}}{\hfill$\square$}
\newenvironment{solution3}{\paragraph{\emph{Solution $3$}.}}{\hfill$\square$}

\begin{document} 

\title[Homework 5]{Homework 5}

\date{\today} 
\maketitle 


The following problem is Theorem 2.27 in your book.  Try to prove it first on your own before looking at the book.
\begin{problem}
Suppose $X$ is a set and $\mathcal{S}$ is a set of subsets of $X$ (that is $\mathcal{S} \subseteq \mathcal{P}(X)$).  Show that the intersection of all $\sigma$-algebras on $X$ that contain $\mathcal{S}$ is a $\sigma$-algebra on $X$.
\end{problem}
\begin{solution1}

\end{solution1}

The $\sigma$-algebra from the last problem is called the $\sigma$-algebra generated by $\mathcal{S}$; it is the smallest (with respect to inclusion) $\sigma$-algebra on $X$ that contains $\mathcal{S}$.  In particular, the smallest $\sigma$-algebra on $\mathbb{R}^{n}$ containing all the open sets of $\mathbb{R}^{n}$ is called the Borel $\sigma$-algebra on $\mathbb{R}^{n}$ and is denoted by $\mathcal{B}(\mathbb{R}^{n})$.  Sets in $\mathcal{B}(\mathbb{R}^{n})$ are called Borel sets.  Since we showed in class that every open set is Lebesgue measurable, we have the inclusion
$$\mathcal{B}(\mathbb{R}^{n}) \subseteq \mathcal{L}(\mathbb{R}^{n}) $$
of $\sigma$-algebras.  An immediate question that arises is:  are these two $\sigma$-algebras equal?  (Spoiler alert:  no.  There are sets that are Lebesgue measurable, but are not Borel sets.  In a few weeks, you should be able to write down such a set.)
\begin{problem}
Is the Cantor set $C \subseteq [0,1]$ a Borel set of $\mathbb{R}$??  Explain carefully.
\end{problem}
\begin{solution1}

\end{solution1}


\begin{problem}[2B-3]
Suppose $\mathcal{A}$ is the smallest $\sigma$-algebra on $\mathbb{R}$ containing $\{(r,s] : r,s \in \mathbb{Q} \}$.  Prove that $\mathcal{A} = \mathcal{B}(\mathbb{R})$.  (You can also attempt problems 4,5,6 of \S2B.)
\end{problem}
\begin{solution1}

\end{solution1}

In class, we defined the notion of a measure space $(X,\mathcal{A},\mu)$.  We can think of this as an abstraction of $(\mathbb{R}^{n},\mathcal{L}(\mathbb{R}^{n}),\lambda)$.  The following problem is Theorem 2.57 in your book, but go ahead and try to prove it yourself first.
\begin{problem}
Let $(X,\mathcal{A},\mu)$ be a measure space and let $A,B \in \mathcal{A}$ be such that $A \subseteq B$.  Show that
\begin{enumerate}
\item $\mu(A) \le \mu(B) $
\item $\mu(B \smallsetminus A) = \mu(B) - \mu(A)$ provided that $\mu(A) < \infty$.
\end{enumerate}
\end{problem}
\begin{solution1}

\end{solution1}

The following problem is Theorem 2.58 in your book, but go ahead and try to prove it yourself first.
\begin{problem}
Let $(X,\mathcal{A},\mu)$ be a measure space and let $A_{1}, A_{2}, \ldots \in \mathcal{A}$.  Show that
$$\mu \left( \bigcup_{k=1}^{\infty}A_{k}\right) \le \sum_{k=1}^{\infty} \mu(A_{k}). $$
(In other words, an abstract measure is alway countably subadditive.)
\end{problem}
\begin{solution1}

\end{solution1} 

Similarly, the following problem is Theorem 2.59 in your book, but go ahead and try to prove it yourself first.
\begin{problem}
Let $(X,\mathcal{A},\mu)$ be a measure space and let
$$A_{1} \subseteq A_{2} \subseteq \ldots $$
be an increasing sequence of sets in $\mathcal{A}$.  Show that
$$\mu \left( \bigcup_{k=1}^{\infty}A_{k}\right) = \lim_{k \to \infty} \mu(A_{k}). $$

\end{problem}
\begin{solution1}

\end{solution1}


Once more, the following problem is Theorem 2.60 in your book, but go ahead and try to prove it yourself first.
\begin{problem}
Let $(X,\mathcal{A},\mu)$ be a measure space and let
$$A_{1} \supseteq A_{2} \supseteq \ldots $$
be a decreasing sequence of sets in $\mathcal{A}$ with $\mu(A_{1}) < \infty$.  Then
$$\mu \left( \bigcap_{k=1}^{\infty} A_{k}\right) = \lim_{k \to \infty} \mu(A_{k}). $$
\end{problem}
\begin{solution1}

\end{solution1}

\begin{problem}[2C-10]
Give an example of a measure space $(X,\mathcal{A},\mu)$ and a decreasing sequence
$$A_{1} \supseteq A_{2} \supseteq \ldots $$
of sets in $\mathcal{A}$ such that
$$\mu \left( \bigcap_{k=1}^{\infty} A_{k}\right) \neq \lim_{k \to \infty} \mu(A_{k}). $$
\end{problem}
\begin{solution1}

\end{solution1}

By now you could read and understand Sections 1A, 1B, 2A, 2B except the section on measurable functions, 2C and 2D and compare Axler's presentation with what we did in class so far.  

\end{document}









